\chapter{Valoriser une option : d'abord le modèle, ensuite la méthode}

\section{Principe commun : l'espérance risque-neutre}

Évaluer une option, c'est calculer l'\textbf{espérance actualisée de son payoff sous la probabilité risque-neutre} $\Q$ :
\[
V_0 = e^{-rT}\,\E_\Q[\text{Payoff}].
\]
Deux ingrédients, à ne pas confondre, interviennent \emph{dans cet ordre} :
\begin{enumerate}[nosep]
\item \textbf{le modèle} dit \emph{comment évolue le prix du sous-jacent} (volatilité constante ou aléatoire, sauts\dots) ; il fixe la loi de $S_T$ ou de la trajectoire, donc ce que vaut l'espérance ;
\item \textbf{la méthode} dit \emph{comment on calcule numériquement cette espérance} une fois le modèle choisi (formule, EDP, arbre, simulation\dots).
\end{enumerate}
On choisit donc toujours le modèle d'abord : sans modèle, il n'y a pas d'espérance à calculer. Une même méthode (Monte-Carlo) peut servir pour tous les modèles, et un même modèle (Black-Scholes) peut être traité par toutes les méthodes.

\section{Les modèles du sous-jacent}

\begin{table}[H]
\centering\small
\begin{tabularx}{\textwidth}{>{\bfseries}p{2.7cm}p{6.2cm}L}
\toprule
Modèle & Dynamique sous $\Q$ & Prix du call européen \\
\midrule
Black-Scholes (GBM) & $dS_t = rS_tdt + \sigma S_tdW_t$ & Formule fermée \\
Volatilité locale (Dupire) & $dS_t = rS_tdt + \sigma(t,S_t)S_tdW_t$ & EDP ou Monte-Carlo \\
Heston & $dS_t = rS_tdt + \sqrt{v_t}S_tdW^1_t$,\ \ $dv_t = \kappa(\theta-v_t)dt + \xi\sqrt{v_t}dW^2_t$ & Semi-fermée (Fourier) \\
Merton (sauts) & $\frac{dS_t}{S_{t^-}} = (r-\lambda\kappa)dt + \sigma dW_t + (J-1)dN_t$ & Série de prix BS \\
Vasicek / CIR & taux court $dr_t = a(b-r_t)dt + \sigma_r r_t^{\gamma}dW_t$ & Hors cadre (options sur taux) \\
\bottomrule
\end{tabularx}
\caption{Modèles de dynamique du sous-jacent et méthodes disponibles.}
\label{tab:modeles}
\end{table}

\begin{itemize}[nosep]
\item \textbf{Black-Scholes} : volatilité constante, trajectoires continues. C'est le modèle le plus simple et le seul qui donne une formule fermée pour l'européenne. \emph{C'est notre cadre.}
\item \textbf{Dupire} : la volatilité dépend du temps et du prix, ce qui permet de reproduire exactement les prix d'options observés sur le marché.
\item \textbf{Heston} : la variance $v_t$ est elle-même aléatoire (retour à la moyenne $\theta$ à la vitesse $\kappa$, volatilité de la volatilité $\xi$) ; elle explique le « smile » de volatilité.
\item \textbf{Merton} : ajoute des sauts (processus de Poisson $N_t$ d'intensité $\lambda$, amplitude $J$) pour représenter les krachs.
\item \textbf{Vasicek / CIR} : modèles de \emph{taux} d'intérêt, hors de notre cadre (taux $r$ constant).
\end{itemize}
Quel que soit le modèle, Monte-Carlo s'applique de la même façon : seule la simulation de $S_T$ change. Nous n'implémentons que Black-Scholes, que nous détaillons maintenant.

\section{Le modèle de Black-Scholes : le prix exact de référence}

\subsection{L'équation différentielle stochastique}
Sous la mesure risque-neutre, le prix du sous-jacent suit le mouvement brownien géométrique
\[
dS_t = r\,S_t\,dt + \sigma\,S_t\,dW_t ,\qquad S_0 \text{ donné},
\]
où $W$ est un mouvement brownien : $W_0=0$, trajectoires continues, accroissements indépendants et $W_t-W_s\sim\Ncal(0,t-s)$. Le terme $rS_tdt$ est la tendance (taux sans risque) et $\sigma S_tdW_t$ le bruit (volatilité $\sigma$).

\subsection{Existence et unicité}
Les coefficients $a(x) = rx$ et $b(x) = \sigma x$ sont lipschitziens et à croissance linéaire :
$|a(x)-a(y)| \le r|x-y|$, $|b(x)-b(y)|\le\sigma|x-y|$, $|a(x)|+|b(x)|\le (r+\sigma)(1+|x|)$.
Le théorème d'existence et d'unicité des EDS garantit donc une solution forte unique, continue, pour tout $S_0$.
\emph{Probabilité réelle ou risque-neutre ?} Sous la probabilité réelle la tendance est $\mu$ ; pour valoriser, le théorème de Girsanov permet de passer à $\Q$ sous laquelle la tendance devient $r$.

\subsection{Résolution : la loi log-normale}
\begin{demo}[ — solution de l'EDS par le lemme d'Itô]
On pose $X_t = \ln S_t$, soit $f(s)=\ln s$, $f'(s) = 1/s$, $f''(s) = -1/s^2$. Le lemme d'Itô donne
\[
dX_t = f'(S_t)dS_t + \tfrac12f''(S_t)(dS_t)^2 = \frac{dS_t}{S_t} - \frac{(dS_t)^2}{2S_t^2}.
\]
Comme $(dW_t)^2 = dt$ et $dt\,dW_t = (dt)^2 = 0$, on a $(dS_t)^2 = \sigma^2S_t^2dt$, d'où
\[
dX_t = \Big(r-\frac{\sigma^2}{2}\Big)dt + \sigma dW_t.
\]
Les coefficients étant constants, on intègre entre $0$ et $T$, puis on écrit $W_T = \sqrt T Z$ avec $Z\sim\Ncal(0,1)$ :
\[
\boxed{S_T = S_0\exp\Big[\Big(r-\frac{\sigma^2}2\Big)T + \sigma\sqrt T\,Z\Big]}
\qquad\Longrightarrow\qquad \ln S_T\sim\Ncal\Big(\ln S_0+\Big(r-\frac{\sigma^2}2\Big)T,\ \sigma^2T\Big).
\]
\end{demo}

$S_T$ est donc \textbf{log-normale}. Le terme $-\sigma^2/2$ est la \emph{correction d'Itô} : comme $\E[e^{aZ}] = e^{a^2/2}$, on obtient
$\E[S_T] = S_0e^{(r-\sigma^2/2)T}e^{\sigma^2T/2} = S_0e^{rT}$. L'actif actualisé $e^{-rt}S_t$ est d'espérance constante (c'est une martingale) : c'est la propriété qui définit la mesure risque-neutre.

\subsection{La formule de Black-Scholes}
\begin{demo}[ — formule du call européen]
On veut $C = e^{-rT}\E[(S_T-K)^+]$.

\emph{a) Événement d'exercice.} $S_T>K \iff Z>-d_2$, avec $d_2 = \dfrac{\ln(S_0/K)+(r-\sigma^2/2)T}{\sigma\sqrt T}$. Donc $\Q(S_T>K) = N(d_2)$.

\emph{b) Découpage.} $\E[(S_T-K)^+] = \E[S_T\1_{\{Z>-d_2\}}] - K\,N(d_2)$.

\emph{c) Premier terme : on complète le carré.} En utilisant $\sigma\sqrt Tz - \frac{z^2}{2} = -\frac12(z-\sigma\sqrt T)^2 + \frac{\sigma^2T}{2}$ :
\[
\E[S_T\1_{\{Z>-d_2\}}] = S_0e^{(r-\frac{\sigma^2}{2})T}\!\int_{-d_2}^{\infty}\! e^{\sigma\sqrt Tz}\frac{e^{-z^2/2}}{\sqrt{2\pi}}dz
= S_0e^{rT}\!\int_{-d_2}^{\infty}\!\frac{e^{-(z-\sigma\sqrt T)^2/2}}{\sqrt{2\pi}}dz = S_0e^{rT}N(d_2+\sigma\sqrt T).
\]
\emph{d) Conclusion.} Avec $d_1 = d_2+\sigma\sqrt T$ :
\[
\boxed{\CBS = S_0N(d_1) - Ke^{-rT}N(d_2)},\qquad d_1 = \frac{\ln(S_0/K)+(r+\sigma^2/2)T}{\sigma\sqrt T},\quad d_2 = d_1-\sigma\sqrt T.
\]
\end{demo}

$N(\cdot)$ est la fonction de répartition de $\Ncal(0,1)$ et $N(d_2)$ la probabilité risque-neutre de finir dans la monnaie.

\textbf{Exemple chiffré (référence de tout le projet).} $d_1 = (0 + 0{,}07)/0{,}2 = 0{,}35$, $d_2 = 0{,}15$, $N(d_1)\approx 0{,}6368$, $N(d_2)\approx0{,}5596$ :
\[
\CBS = 100\times0{,}6368 - 100\,e^{-0{,}05}\times0{,}5596 = \mathbf{10{,}4506}.
\]


\subsection{L'EDP de Black-Scholes et la formule de Feynman-Kac}
Le prix $V(t,S)$ vérifie aussi une équation aux dérivées partielles. En appliquant Itô à $e^{-rt}V(t,S_t)$, qui doit être une martingale sous $\Q$, le terme en $dt$ doit s'annuler :
\[
\frac{\partial V}{\partial t} + \frac12\sigma^2S^2\frac{\partial^2V}{\partial S^2} + rS\frac{\partial V}{\partial S} - rV = 0,\qquad V(T,S) = (S-K)^+.
\]
Réciproquement, la \textbf{formule de Feynman-Kac} dit que la solution de cette EDP s'écrit comme une espérance :
$V(t,S) = e^{-r(T-t)}\E_\Q[(S_T-K)^+\mid S_t = S]$.
\emph{Résoudre l'EDP (analyse) et calculer l'espérance (probabilités) sont donc deux façons de résoudre le même problème} : c'est ce qui justifie les différentes méthodes ci-dessous.

\section{Les méthodes de valorisation}

Le modèle étant fixé, il reste à calculer l'espérance. Nous détaillons chaque méthode : principe, mise en œuvre, application au call européen, limites.

\subsection{Formule fermée}
\textbf{Principe.} Calculer l'intégrale exactement, comme dans la démonstration ci-dessus.
\textbf{Call européen :} oui, c'est $\CBS$, instantanée et exacte.
\textbf{Limite :} n'existe que dans des cas très particuliers (modèle simple, payoff qui ne dépend que de $S_T$). Dès que le payoff dépend de toute la trajectoire, il n'y a en général plus de formule.

\subsection{EDP et différences finies}
\textbf{Principe.} On résout numériquement l'EDP de Black-Scholes sur une grille $S_i = i\Delta S$ ($0\le i\le n_S$, $S_{\max} = 2K$), $t_n = n\Delta t$ ($0\le n\le n_t$), en partant de la condition finale $V(T,S) = (S-K)^+$ et en remontant le temps. Les dérivées sont remplacées par des différences :
\[
\frac{\partial V}{\partial S}\approx\frac{V_{i+1}-V_{i-1}}{2\Delta S},\qquad \frac{\partial^2V}{\partial S^2}\approx\frac{V_{i+1}-2V_i+V_{i-1}}{\Delta S^2}.
\]
\textbf{Schéma implicite.} Ces dérivées sont prises au temps inconnu $t_n$, ce qui donne à chaque pas un système tridiagonal
\[
a_iV^n_{i-1} + b_iV^n_i + c_iV^n_{i+1} = V^{n+1}_i,\quad a_i = \tfrac{\Delta t}{2}(ri - \sigma^2i^2),\ b_i = 1+\Delta t(\sigma^2i^2+r),\ c_i = -\tfrac{\Delta t}{2}(\sigma^2i^2+ri),
\]
avec les conditions aux bords $V(t,0)=0$ et $V(t,S_{\max}) = S_{\max}-Ke^{-r(T-t)}$. On le résout en $O(n_S)$ opérations par l'algorithme de Thomas. Le schéma implicite est inconditionnellement stable (contrairement au schéma explicite, qui impose $\Delta t \lesssim \Delta S^2/(\sigma^2S_{\max}^2)$).


\textbf{Atout :} l'erreur diminue quand on affine la grille, et en dimension 1 l'EDP est \emph{plus efficace} que Monte-Carlo.
\textbf{Limite :} \emph{malédiction de la dimension} : avec $d$ variables d'état (plusieurs sous-jacents, volatilité aléatoire, payoff dépendant du chemin\dots) et $n$ points par axe, la grille a $n^d$ nœuds.

\subsection{Arbre binomial de Cox-Ross-Rubinstein}
\textbf{Principe.} On découpe $[0,T]$ en $n$ pas. À chaque pas, le prix est multiplié par $u = e^{\sigma\sqrt{\Delta t}}$ ou par $d = 1/u$. La probabilité de hausse
\[
p = \frac{e^{r\Delta t}-d}{u-d}
\]
est choisie pour que $\E[S_{t+\Delta t}\mid S_t] = S_te^{r\Delta t}$ (neutralité au risque). On calcule les payoffs aux $n+1$ feuilles, puis on remonte : $V = e^{-r\Delta t}\big(pV_{\text{haut}}+(1-p)V_{\text{bas}}\big)$. Quand $n\to\infty$, la loi de $S_T$ dans l'arbre tend vers la log-normale (théorème central limite) et le prix tend vers $\CBS$.


\textbf{Atouts :} très intuitif, et c'est la méthode naturelle pour les options \emph{américaines} (à chaque nœud on compare la valeur de continuation à la valeur d'exercice).
\textbf{Limite :} convergence lente et oscillante en $O(1/n)$ ; et si le payoff dépend du chemin, l'arbre ne se « recombine » plus et il faut suivre les $2^n$ chemins.

\subsection{Monte-Carlo}
\textbf{Principe.} Remplacer l'espérance par une moyenne de payoffs simulés (loi des grands nombres), et mesurer l'erreur par un intervalle de confiance (théorème central limite). La méthode est détaillée au chapitre 3.
\textbf{Call européen :} oui ; ici la formule exacte permet de mesurer précisément son erreur.
\textbf{Atout :} l'erreur reste en $\sigma_X/\sqrt N$ \emph{quelle que soit la dimension} du problème, et la méthode s'adapte à n'importe quel modèle ou payoff.
\textbf{Limite :} convergence lente en $1/\sqrt N$, d'où la réduction de variance.

\subsection{Transformée de Fourier (Carr-Madan)}
\textbf{Principe.} Si l'on connaît la fonction caractéristique $\varphi(u) = \E[e^{iu\ln S_T}]$ (c'est le cas pour Black-Scholes, Heston, Merton), le prix d'un call amorti $e^{\alpha k}C(k)$, $k = \ln K$, a une transformée de Fourier explicite :
\[
\psi(v) = \frac{e^{-rT}\varphi\big(v-(\alpha+1)i\big)}{\alpha^2+\alpha-v^2+i(2\alpha+1)v},\qquad C(k) = \frac{e^{-\alpha k}}{\pi}\int_0^\infty \Re\big(e^{-ivk}\psi(v)\big)\,dv,
\]
calculée pour tous les strikes en une seule FFT.
\textbf{Call européen :} oui, si $\varphi$ est connue ; très rapide pour calibrer un modèle sur de nombreux strikes.
\textbf{Limite :} exige une fonction caractéristique connue, et ne s'applique qu'aux payoffs qui dépendent de $S_T$ seulement.

\subsection{Synthèse comparative}

\begin{table}[H]
\centering\small
\begin{tabularx}{\textwidth}{>{\bfseries}p{2.6cm}LLL}
\toprule
Méthode & Call européen & Atout & Limite principale \\
\midrule
Formule fermée & Oui : Black-Scholes, exacte & Instantanée & Cas très particuliers \\
EDP (différences finies) & Oui, converge vers $\CBS$ & Très efficace en dimension 1 & Malédiction de la dimension \\
Arbre CRR & Oui, converge en $O(1/n)$ & Intuitif, options américaines & Lent, oscillant en $O(1/n)$ \\
Monte-Carlo & Oui, IC $\pm0{,}029$ ($N=10^6$) & Indépendant de la dimension & Erreur en $1/\sqrt N$ \\
Fourier (Carr-Madan) & Oui si $\varphi$ connue & Tous les strikes en une FFT & Exige $\varphi$ \\
\bottomrule
\end{tabularx}
\caption{Synthèse des méthodes de valorisation appliquées au call européen.}
\end{table}

\begin{retenir}
Pour le call européen, toutes les méthodes fonctionnent et la formule fermée est imbattable : elle sert de \emph{référence}. Monte-Carlo est la plus lente ici, mais c'est la seule dont l'erreur ne dépend pas de la dimension du problème. D'où notre démarche : \textbf{valider et optimiser Monte-Carlo sur le call européen}, là où l'on connaît la bonne réponse.
\end{retenir}
